\chapter{Ausarbeitung des linearen Problems}
\label{kap:modell}

Kapitel~2 hat die Werkzeuge bereitgestellt, mit denen sich ein
Zuordnungsproblem als ganzzahliges lineares Programm fassen lässt. Das
vorliegende Kapitel wendet diese Werkzeuge auf die Fragestellung der Erpolino
Metallverarbeitung GmbH an. Es entwickelt zunächst die Annahmen, unter denen die
betriebliche Situation überhaupt einer mathematischen Beschreibung zugänglich
wird, formuliert anschließend Mengen, Parameter und Variablen, stellt
Zielfunktion und Nebenbedingungen auf und diskutiert die Eigenschaften der
gewonnenen Formulierung. Den Abschluss bildet eine methodische Gegenüberstellung
der Verfahren, die zur Lösung des Modells in Frage kommen.

\section{Modellannahmen}
\label{sec:annahmen}

Ein Modell ist niemals ein Abbild der Wirklichkeit, sondern eine bewusste
Vereinfachung. Welche Vereinfachungen zulässig sind, entscheidet sich daran, ob
sie den Kern der Fragestellung berühren. Für die hier untersuchte Zuordnung von
Arbeitsgängen zu Mitarbeitern werden fünf Annahmen getroffen, die im Folgenden
begründet werden.

Erstens bleiben \Index{Rüstzeiten} unberücksichtigt. Die Bearbeitungsdauer eines
Arbeitsgangs hängt allein davon ab, welche Person ihn ausführt, nicht davon,
welche Tätigkeit dieser Person unmittelbar vorausging. In der Fertigung der
Erpolino Metallverarbeitung GmbH fallen Rüstvorgänge zwar an, sie sind jedoch
überwiegend an das Betriebsmittel und nicht an die Reihenfolge der Arbeitsgänge
gebunden und daher bereits in den erhobenen Grundzeiten enthalten. Wären
reihenfolgeabhängige Rüstzeiten zu berücksichtigen, so entstünde ein Problem mit
Rundreisecharakter, dessen Komplexität die des Zuordnungsproblems deutlich
übersteigt \parencite{pinedo2016scheduling}.

Zweitens werden keine \Index{Reihenfolgeabhängigkeiten} modelliert. Alle
Arbeitsgänge gelten als voneinander unabhängig und zu jedem Zeitpunkt ausführbar.
Diese Annahme ist der eigentliche Grund dafür, dass das Problem ein reines
Zuordnungsproblem bleibt und keine zeitliche Einplanung erfordert. Sobald
feststeht, welche Person welche Arbeitsgänge übernimmt, ist die Reihenfolge
innerhalb einer Person für die Zielgröße bedeutungslos, denn die Summe der
Bearbeitungszeiten ändert sich durch Umsortieren nicht. Damit entfallen die
Zeitpunktvariablen, die ein Ablaufplanungsmodell sonst mit sich brächte.

Drittens sind die Aufgaben \Index{unteilbar}. Ein Arbeitsgang wird vollständig
von genau einer Person ausgeführt; eine Aufteilung auf zwei Mitarbeiter oder ein
Wechsel während der Bearbeitung ist ausgeschlossen. Diese Annahme entspricht der
betrieblichen Praxis, in der ein Arbeitsgang zusammen mit dem zugehörigen
Werkstück und dem Prüfprotokoll einer Person zugewiesen wird. Sie ist zugleich
die Ursache dafür, dass das Modell ganzzahlige Variablen benötigt und nicht als
lineares Programm gelöst werden kann. Ließe man Teilbarkeit zu, so wäre das
Problem in polynomieller Zeit lösbar; die \Index{Ganzzahligkeitsforderung} ist
also nicht Beiwerk, sondern der Grund für die Schwierigkeit des Problems
\parencite{garey1979computers}.

Viertens stehen alle Mitarbeiter über den Planungshorizont hinweg gleichermaßen
zur Verfügung. Schichtmodelle, Pausen, Urlaub und Krankheit bleiben außer
Betracht. Der Planungshorizont ist eine Kalenderwoche, für die die
Arbeitsvorbereitung ein Fertigungslos freigibt; innerhalb dieser Woche gilt jede
Person als ganztags einsatzbereit.

Fünftens ist die \Index{Effizienz} einer Person je Qualifikation konstant.
Lerneffekte, Ermüdung und Tagesform werden nicht abgebildet. Die
Bearbeitungszeit ergibt sich damit deterministisch aus der Grundzeit des
Arbeitsgangs und einem personenbezogenen Effizienzfaktor. Stochastische
Schwankungen der Bearbeitungsdauer wären grundsätzlich modellierbar, führten
aber auf ein stochastisches Optimierungsproblem und damit auf eine andere
Klasse von Verfahren.

Diese Annahmen idealisieren die Fertigung erheblich. Sie lassen jedoch genau
jenen Effekt unangetastet, um den es der Fertigungsleitung geht: die ungleiche
Auslastung, die entsteht, wenn Qualifikationen ungleich verteilt sind.

\section{Mengen, Parameter und Variablen}
\label{sec:mengen}

Sei $I$ die Menge der Mitarbeiter der Fertigung und $J$ die Menge der
Arbeitsgänge des Planungshorizonts. Für die untersuchte Instanz gilt
$|I| = \ZahlMitarbeiter{}$ und $|J| = \ZahlAufgaben{}$. Jeder Arbeitsgang
$j \in J$ verlangt genau eine der \ZahlQualifikationen{} Qualifikationen; jeder
Mitarbeiter $i \in I$ beherrscht eine Teilmenge davon.

Aus dieser Qualifikationsstruktur ergibt sich die zentrale Menge des Modells. Es
bezeichne
\begin{equation}
  Z \subseteq I \times J
  \label{eq:zulaessig}
\end{equation}
die Menge der \Index{zulässigen Paare}, also derjenigen Kombinationen $(i,j)$,
für die Mitarbeiter $i$ die von Arbeitsgang $j$ geforderte Qualifikation besitzt.
In der Scheduling-Notation entspricht dies der Restriktion $M_j$, die für jeden
Auftrag die Menge der zulässigen Maschinen einschränkt
\parencite{pinedo2016scheduling}. Die Menge $Z$ ist keine bloße
Schreibvereinfachung, sondern trägt die gesamte Information über die
\Index{Qualifikationsschranken}; ihre Rolle wird in
Abschnitt~\ref{sec:eigenschaften} ausführlich behandelt.

Für jedes zulässige Paar $(i,j) \in Z$ ist die \Index{Bearbeitungszeit}
$p_{ij} > 0$ in Minuten definiert. Sie wird nicht frei erhoben, sondern aus zwei
Größen berechnet: der Grundzeit $g_j$ des Arbeitsgangs, die die
Arbeitsvorbereitung für eine Normalleistung vorgibt, und dem Effizienzfaktor
$e_i(q_j)$, mit dem Mitarbeiter $i$ die von Arbeitsgang $j$ geforderte
Qualifikation $q_j$ ausübt. Es gilt
\begin{equation}
  p_{ij} = \frac{g_j}{e_i(q_j)}
  \qquad \text{für alle } (i,j) \in Z .
  \label{eq:bearbeitungszeit}
\end{equation}
Ein Effizienzfaktor von $1{,}00$ bedeutet Normalleistung, ein Wert über $1{,}00$
eine überdurchschnittliche und ein Wert darunter eine unterdurchschnittliche
Leistung. Ein Mitarbeiter mit dem Faktor $1{,}25$ benötigt für einen
Arbeitsgang mit der Grundzeit $95$ Minuten folglich $76$ Minuten. Wichtig ist,
dass die Effizienzfaktoren nicht in einer festen Rangfolge stehen: Eine Person
kann beim Drehen überdurchschnittlich und beim Schleifen unterdurchschnittlich
sein. Genau deshalb liegt kein Problem mit gleichförmigen Maschinen vor, bei dem
sich alle Bearbeitungszeiten aus einer einzigen personenbezogenen Geschwindigkeit
ergäben, sondern ein Problem mit \Index{unabhängigen parallelen Maschinen}, in
der Kendall-artigen Notation $R \mid M_j \mid C_{\max}$
\parencite{pinedo2016scheduling}.

Das Modell benötigt zwei Arten von Variablen. Zum einen die
\Index{Binärvariable}
\begin{equation}
  x_{ij} \in \{0,1\}
  \qquad \text{für alle } (i,j) \in Z ,
  \label{eq:xvar}
\end{equation}
die den Wert $1$ annimmt, wenn Mitarbeiter $i$ den Arbeitsgang $j$ übernimmt, und
sonst $0$. Zum anderen die stetige Variable
\begin{equation}
  C_{\max} \geq 0 ,
  \label{eq:cmaxvar}
\end{equation}
die den \Index{Makespan} abbildet, also den Zeitpunkt, zu dem auch der zuletzt
fertig werdende Mitarbeiter seine Arbeit abgeschlossen hat. Bemerkenswert ist,
dass $C_{\max}$ keine Entscheidungsvariable im eigentlichen Sinne ist: Ihr Wert
ist durch die Belegung der $x_{ij}$ vollständig bestimmt. Sie dient allein dazu,
das Maximum über die Auslastungen, das für sich genommen keine lineare Funktion
ist, in eine lineare Formulierung zu überführen.

\section{Zielfunktion und Nebenbedingungen}
\label{sec:zielfunktion}

Das vollständige Modell lautet
\begin{align}
  \min \quad & C_{\max} \label{eq:ziel} \\
  \text{u.\,d.\,N.} \quad
  & \sum_{i \,:\, (i,j) \in Z} x_{ij} = 1
    && \text{für alle } j \in J , \label{eq:vergabe} \\
  & \sum_{j \,:\, (i,j) \in Z} p_{ij}\, x_{ij} \leq C_{\max}
    && \text{für alle } i \in I , \label{eq:auslastung} \\
  & x_{ij} \in \{0,1\}
    && \text{für alle } (i,j) \in Z , \label{eq:binaer} \\
  & C_{\max} \geq 0 . \label{eq:nichtneg}
\end{align}

Die Zielfunktion~\eqref{eq:ziel} minimiert den Makespan. Sie besteht aus einer
einzigen Variablen und ist damit denkbar einfach; die gesamte Modellierungsarbeit
steckt in den Nebenbedingungen.

Die \Index{Vergabebedingung}~\eqref{eq:vergabe} stellt sicher, dass jeder
Arbeitsgang genau einmal vergeben wird. Die Summation läuft ausschließlich über
diejenigen Mitarbeiter, für die das Paar $(i,j)$ zulässig ist. Die Gleichung
verlangt zweierlei zugleich: Kein Arbeitsgang bleibt liegen, und kein
Arbeitsgang wird doppelt ausgeführt. Da die Menge $Z$ zu jedem $j \in J$
mindestens einen Mitarbeiter enthält -- jede Qualifikation ist in der Fertigung
mehrfach vorhanden --, ist das Modell zulässig.

Die \Index{Auslastungsbedingung}~\eqref{eq:auslastung} verknüpft die
Zuordnungsvariablen mit dem Makespan. Die linke Seite summiert die
Bearbeitungszeiten aller Arbeitsgänge, die Mitarbeiter $i$ übernimmt, und
beschreibt damit dessen Auslastung. Für jeden Mitarbeiter wird gefordert, dass
diese Auslastung den Wert $C_{\max}$ nicht übersteigt. Es gibt genau
$\ZahlMitarbeiter{}$ solcher Bedingungen, eine je Mitarbeiter, und
$\ZahlAufgaben{}$ Vergabebedingungen.

Die Formulierung enthält weder Zeitpunkte noch Reihenfolgevariablen. Sie ist
damit erheblich kleiner als ein allgemeines Ablaufplanungsmodell und verdankt
diese Kompaktheit unmittelbar den Annahmen aus Abschnitt~\ref{sec:annahmen}.

\section{Eigenschaften des Modells}
\label{sec:eigenschaften}

Vier Eigenschaften der Formulierung verdienen eine genauere Betrachtung, weil sie
weniger selbstverständlich sind, als sie auf den ersten Blick erscheinen.

\subsection*{Die Ungleichung in der Auslastungsbedingung}

Warum steht in~\eqref{eq:auslastung} ein Kleiner-gleich-Zeichen und kein
Gleichheitszeichen? Die Frage ist berechtigt, denn $C_{\max}$ soll ja gerade das
Maximum der Auslastungen sein, und ein Maximum wird von mindestens einem
Mitarbeiter angenommen.

Der Punkt ist, dass die Ungleichung zusammen mit der Minimierung genau das
Gewünschte leistet. Die Bedingungen~\eqref{eq:auslastung} erzwingen, dass
$C_{\max}$ mindestens so groß ist wie jede einzelne Auslastung, also
$C_{\max} \geq \max_{i \in I} \sum_{j} p_{ij} x_{ij}$. Die Zielfunktion drückt
$C_{\max}$ so weit nach unten, wie es die Bedingungen erlauben. In jeder
Optimallösung gilt daher Gleichheit für mindestens einen Mitarbeiter, und
$C_{\max}$ nimmt exakt den Wert des Maximums an. Die Kombination aus unterer
Schranke und Minimierungsrichtung ist die übliche lineare Darstellung einer
Maximumsfunktion \parencite{wolsey2020integer}.

Ein Gleichheitszeichen wäre dagegen schlicht falsch. Es würde verlangen, dass
alle \ZahlMitarbeiter{} Mitarbeiter exakt dieselbe Auslastung erreichen. Bei
unteilbaren Aufgaben und personenabhängigen Bearbeitungszeiten ist eine solche
perfekt gleichmäßige Verteilung im Allgemeinen nicht erreichbar; das Modell wäre
unzulässig und der Solver meldete, dass keine Lösung existiert -- obwohl die
Fertigung selbstverständlich jeden Tag zulässige Pläne umsetzt. Der Unterschied
zwischen $\leq$ und $=$ entscheidet hier also nicht über Eleganz, sondern über
Lösbarkeit.

\subsection*{Unzulässige Paare existieren nicht}

Eine naheliegende Alternative zur Menge $Z$ bestünde darin, für alle
$\ZahlMitarbeiter{} \times \ZahlAufgaben{} = 192$ Kombinationen eine Variable
anzulegen und unzulässige Zuordnungen durch eine sehr große Bearbeitungszeit zu
bestrafen. Dieses \Index{Big-M}-Vorgehen ist verbreitet und in vielen
Modellierungssituationen unvermeidlich -- hier ist es jedoch ein Nachteil.

Der Grund liegt in der \Index{LP-Relaxation}. Verfahren vom Typ
\Index{Branch-and-Bound} lösen zunächst das Modell ohne
Ganzzahligkeitsforderung, verwenden dessen Zielfunktionswert als untere Schranke
und verzweigen anschließend über fraktionale Variablen
\parencite{land1960automatic,nemhauser1988integer}. Die Qualität dieser Schranke
bestimmt maßgeblich, wie viele Knoten der Suchbaum umfasst. Eine große Konstante
$M$ erzeugt nun eine Relaxation, in der die bestrafte Variable einen kleinen
Bruchteil ihres Wertes annehmen kann, ohne die Zielfunktion nennenswert zu
belasten. Der relaxierte Zielfunktionswert liegt dann weit unter dem
ganzzahligen Optimum, die \Index{Dualitätslücke} ist groß, und der Suchbaum
wächst entsprechend. Big-M-Formulierungen sind für schwache Relaxationen
bekannt \parencite{wolsey2020integer}.

Lässt man die unzulässigen Paare dagegen von vornherein weg, so kann die
Relaxation gar nicht erst auf sie ausweichen. Die Schranke wird schärfer, weil
sie über einem kleineren und wahrheitsgetreuen Lösungsraum gebildet wird.
Zugleich schrumpft das Modell. Der Vorteil ist damit doppelt und in beiden
Hinsichten echt: Das Modell ist kleiner \emph{und} seine Relaxation ist besser.
Dass sich die Qualifikationsschranken so sauber ausdrücken lassen, ist ein
Glücksfall der vorliegenden Struktur; er sollte genutzt werden.

\subsection*{Modellgröße}

Die Auswirkung lässt sich beziffern. Bei \ZahlMitarbeiter{} Mitarbeitern und
\ZahlAufgaben{} Arbeitsgängen wären $192$ Kombinationen denkbar. Da jeder
Mitarbeiter aber nur zwei bis drei der \ZahlQualifikationen{} Qualifikationen
beherrscht, bleiben lediglich \ZahlPaare{} zulässige Paare übrig. Das Modell
besitzt also \ZahlPaare{} Binärvariablen und eine stetige Variable, zusammen mit
$\ZahlAufgaben{} + \ZahlMitarbeiter{} = 32$ Nebenbedingungen. Die
Qualifikationsschranken dünnen das Modell auf weniger als die Hälfte seines
nominellen Umfangs aus.

Diese Beobachtung hat einen bemerkenswerten Nebeneffekt. Die
Qualifikationsschranken erschweren die Planung aus betrieblicher Sicht, weil sie
Handlungsspielraum nehmen. Aus Sicht des Solvers hingegen sind sie eine Hilfe:
Sie verkleinern den Suchraum. Was den Menschen behindert, entlastet die
Maschine.

\subsection*{Untere Schranken und ihre Schärfe}

Eine untere Schranke für den Makespan lässt sich ohne jede Optimierung angeben.
Kein Arbeitsgang kann schneller erledigt werden als von der dafür am besten
geeigneten Person; die dabei anfallende Arbeit muss auf \ZahlMitarbeiter{}
Mitarbeiter verteilt werden. Damit gilt
\begin{equation}
  C_{\max} \;\geq\; \frac{1}{|I|} \sum_{j \in J} \; \min_{i \,:\, (i,j) \in Z} p_{ij}
  \;=\; \UntereSchranke{}\ \text{Minuten}.
  \label{eq:schranke}
\end{equation}
Eine zweite, ebenso einfache Schranke ist die längste unvermeidbare
Einzelaufgabe: Da Aufgaben unteilbar sind, ist der Makespan mindestens so groß
wie
$\max_{j \in J} \min_{i : (i,j) \in Z} p_{ij}$, also so groß wie derjenige
Arbeitsgang, der selbst unter der günstigsten Zuordnung am längsten dauert. Für
die vorliegende Instanz ist die Lastschranke~\eqref{eq:schranke} die stärkere
der beiden; Schranken dieser Art bilden den Ausgangspunkt der klassischen
Gütegarantien für Listenverfahren \parencite{graham1969bounds} und der
Approximationsalgorithmen für unabhängige parallele Maschinen
\parencite{lenstra1990approximation}.

Das tatsächliche Optimum liegt jedoch bei \Optimum{} Minuten. Die Lücke zur
trivialen Schranke beträgt \LueckeSchranke{} Prozent. Die Schranke ist somit
nicht scharf, und dafür gibt es zwei Gründe.

Der erste Grund sind die Qualifikationsschranken. Schranke~\eqref{eq:schranke}
unterstellt, dass jeder Arbeitsgang von seinem jeweils besten Bearbeiter
ausgeführt wird und die Last dennoch gleichmäßig verteilt werden kann. Beides
zugleich ist unmöglich. Weist man jeden Arbeitsgang seinem schnellsten
Bearbeiter zu, so häufen sich die Arbeitsgänge einer Qualifikation bei
derjenigen Person, die diese Qualifikation am besten beherrscht, während andere
Personen leer ausgehen. Wer die Last umverteilen will, muss Arbeitsgänge an
langsamere Personen abgeben und damit die Gesamtarbeit erhöhen. Der Konflikt
zwischen Bestbesetzung und Gleichverteilung ist die eigentliche Schwierigkeit
des Problems -- und genau der Punkt, an dem die Schranke die Wirklichkeit
verfehlt.

Der zweite Grund ist die Unteilbarkeit. Selbst wenn die
Qualifikationsstruktur eine ausgewogene Verteilung zuließe, ließe sich die
mittlere Last~\eqref{eq:schranke} nur dann exakt erreichen, wenn sich die
Arbeitsgänge beliebig fein aufteilen ließen. Da sie das nicht tun, bleibt stets
ein Rest, um den die beste erreichbare Verteilung über dem Mittelwert liegt.

Die Schranke bleibt gleichwohl nützlich. Sie liefert eine Zahl, gegen die sich
jede heuristische Lösung sofort einordnen lässt, und sie ist es, die
Branch-and-Bound überhaupt erst erlaubt, Teilbäume zu verwerfen.

\section{Vergleich nutzbarer Methoden}
\label{sec:verfahrensvergleich}

Für ein Modell dieser Struktur kommen drei Gruppen von Verfahren in Betracht.
Sie unterscheiden sich in Gütegarantie, Rechenzeit, Skalierbarkeit,
Implementierungsaufwand und Erweiterbarkeit.

\Index{Exakte Verfahren} lösen das Modell~\eqref{eq:ziel} bis
\eqref{eq:nichtneg} in seiner ursprünglichen Form. Das Standardverfahren ist
Branch-and-Bound \parencite{land1960automatic}: Es löst die LP-Relaxation,
verzweigt über eine fraktionale Binärvariable und verwirft Teilbäume, deren
Schranke die beste bekannte Lösung nicht mehr unterbieten kann. Der Vorzug ist
die \Index{Gütegarantie}: Das Verfahren liefert nicht nur eine gute Lösung,
sondern zugleich den Beweis, dass keine bessere existiert. Der Preis ist eine
Rechenzeit, die im schlechtesten Fall exponentiell wächst -- das
zugrundeliegende Problem ist \Index{NP-schwer} \parencite{garey1979computers}.
Der Implementierungsaufwand ist gering, sofern ein Solver eingesetzt wird; das
Modell wird deklarativ notiert und von der Software gelöst \parencite{glpk}.
Auch die Erweiterbarkeit ist hoch: Zusätzliche Anforderungen wie Termine oder
Obergrenzen je Person lassen sich als weitere lineare Nebenbedingungen ergänzen,
ohne den Lösungsalgorithmus anzutasten.

\Index{Konstruktionsheuristiken} bauen eine Lösung in einem einzigen Durchlauf
auf und revidieren keine Entscheidung. Das \Index{Greedy}-Verfahren geht die
Arbeitsgänge in gegebener Reihenfolge durch und weist jeden derjenigen
qualifizierten Person zu, die anschließend am frühesten fertig wäre. Die
\ac{lpt}-Regel sortiert die Arbeitsgänge zuvor absteigend nach Dauer und
bearbeitet die langen zuerst; dahinter steht die Einsicht, dass eine spät
eingeplante lange Aufgabe den Makespan unnötig in die Höhe treibt
\parencite{graham1969bounds}. Beide Verfahren sind in wenigen Zeilen
implementiert und laufen in Bruchteilen einer Millisekunde. Eine Gütegarantie
besitzen sie im vorliegenden Fall jedoch nicht: Die klassischen Schranken für
\ac{lpt} gelten für identische Maschinen, während hier unabhängige Maschinen mit
Qualifikationsschranken vorliegen \parencite{lenstra1990approximation}. Ihre
Erweiterbarkeit ist gering -- jede neue Anforderung verlangt einen Eingriff in
die Konstruktionslogik.

\Index{Metaheuristiken} verbessern eine vorhandene Lösung durch wiederholte
lokale Veränderungen. \ac{sa} akzeptiert dabei nicht nur Verbesserungen, sondern
mit einer temperaturabhängigen Wahrscheinlichkeit auch Verschlechterungen, um
lokale Optima wieder verlassen zu können; die Temperatur wird im Verlauf gesenkt,
sodass das Verfahren am Ende nur noch Verbesserungen zulässt
\parencite{kirkpatrick1983optimization}. Verwandte Ansätze sind die
\Index{Tabu-Suche}, die kürzlich besuchte Lösungen für eine gewisse Zahl von
Schritten sperrt und so Zyklen verhindert, sowie \Index{genetische Algorithmen},
die eine Population von Lösungen führen und durch Rekombination und Mutation
weiterentwickeln. Allen gemeinsam ist, dass sie keine Gütegarantie liefern und
über eine Reihe von Parametern verfügen, deren Einstellung Erfahrung verlangt.
Ihre Stärke ist die Skalierbarkeit: Die Rechenzeit wird durch die
Iterationszahl vorgegeben und nicht durch die Instanzgröße bestimmt. Ihre
Erweiterbarkeit ist mittelmäßig; zusätzliche Anforderungen lassen sich in die
Bewertungsfunktion aufnehmen, doch die Nachbarschaft muss weiterhin zulässige
Lösungen erzeugen.

\begin{table}[tb]
  \centering
  % Automatisch erzeugt von Code/tabellen.py -- nicht von Hand ändern.
% Quelle: Code/daten.py und Ergebnisse/*.json

\begin{tabular}{@{}lS[table-format=3.2]S[table-format=1.2]S[table-format=3.1]l@{}}
\toprule
& {Makespan} & {Abstand} & {Rechenzeit} & \\
Verfahren & {min} & {\%} & {ms} & Gütegarantie \\
\midrule
  Untere Schranke (trivial) & 205.44 & -13.02 & {--} & Schranke, keine Lösung \\
  Greedy & 266.67 & 12.90 & 0.0 & keine \\
  LPT & 265.38 & 12.36 & 0.0 & keine \\
  Simulated Annealing & 240.74 & 1.92 & 297.9 & keine \\
  MILP (GLPK) & 236.20 & 0.00 & 676.5 & beweisbar optimal \\
\midrule
\multicolumn{5}{@{}l@{}}{\footnotesize Simulated Annealing über 20 Läufe, je 60\,000 Iterationen:}\\
  \quad Energie nur Makespan & 240.71 & 1.91 & {--} & bestes 239.86, schlecht. 240.84 \\
  \quad Energie mit Glättung & 240.10 & 1.65 & {--} & bestes 236.20, schlecht. 243.68 \\
\bottomrule
\end{tabular}

  \caption[Vergleich der Verfahren]{Die untersuchten Verfahren auf der Instanz von Erpolino}
  \label{tab:verfahren}
\end{table}

Tabelle~\ref{tab:verfahren} stellt die Verfahren mit den Ergebnissen zusammen,
die sie auf der Instanz der Erpolino Metallverarbeitung GmbH erzielen; die
Vorgriffe auf Kapitel~\ref{kap:ergebnisse} sind hier nur als Beleg der
methodischen Argumentation zu lesen. Das Bild ist eindeutig. Das
\ac{milp}-Modell liefert den Makespan \Optimum{} Minuten und damit die
beweisbar optimale Lösung, und zwar in weniger als einer Sekunde. \ac{lpt}
kommt auf \MakespanLPT{} Minuten und verfehlt das Optimum um \AbstandLPT{}
Prozent; \ac{sa} erreicht \MakespanSA{} Minuten und bleibt \AbstandSA{} Prozent
darüber. Die Heuristiken sind schneller, doch bei einer Rechenzeit des exakten
Verfahrens von unter einer Sekunde ist dieser Vorteil ohne praktischen Wert.

Für die vorliegende Problemgröße gibt es damit keinen Grund, auf eine Heuristik
auszuweichen. Die Wahl fällt auf das exakte Verfahren, und zwar nicht, weil es
grundsätzlich überlegen wäre, sondern weil die Instanz klein genug ist, um seine
Schwäche -- die im schlechtesten Fall exponentielle Rechenzeit -- gar nicht erst
wirksam werden zu lassen. Ab welcher Größe sich dieses Verhältnis umkehrt und
die Heuristiken ihre Berechtigung erhalten, untersucht
Kapitel~\ref{kap:ergebnisse} anhand künstlich vergrößerter Instanzen.
